\section{加法、乘法与排列组合}
\label{sec:04-Discrete-Mathematics/01-Combinatorics/01}

给一个模型做网格搜索：优化器三选一，学习率四选一，批大小二选一，跑完所有组合要训练多少次？\(3\times4\times2=24\)。这一步几乎不用想。

改一个条件再问一次。三个优化器里有一个只在两种学习率下收敛，另外两个四种都能用。照旧相乘仍然得 24，实际能跑的却是 \((2+4+4)\times2=20\)。乘法失效的地方不在算术，而在它悄悄用掉的一个前提：第二步有多少个选项，不能依赖第一步选了什么。

计数出错，十有八九不是乘错了数，而是没说清``一个结果''到底长什么样，或者同一个结果被生成了两次。本节把加法、乘法、排列、组合这四件事收回到同一个动作上：先把结果集合和一个更好数的集合对应起来，再数那个好数的。

\subsection{乘法要求分支数整齐，加法要求分支不重叠}

把生成过程画成一棵决策树，条件就看得见了。

\begin{center}
\begin{tikzpicture}[x=1cm,y=1cm,font=\footnotesize]
  \fill (2,3) circle (2pt);
  \foreach \x in {0.6,2,3.4}{\draw[black!55] (2,3) -- (\x,1.7); \fill (\x,1.7) circle (2pt);}
  \foreach \a/\b in {0.6/0.2,0.6/1.0,2/1.6,2/2.4,3.4/3.0,3.4/3.8}{\draw[black!55] (\a,1.7) -- (\b,0.5); \fill (\b,0.5) circle (1.6pt);}
  \node at (2,-0.15) {每个节点分出 2 条：\(3\times2=6\)};
  \fill (8,3) circle (2pt);
  \foreach \x in {6.6,8,9.6}{\draw[black!55] (8,3) -- (\x,1.7); \fill (\x,1.7) circle (2pt);}
  \foreach \a/\b in {6.6/6.2,6.6/7.0,8/8.0,9.6/8.9,9.6/9.7,9.6/10.5}{\draw[black!55] (\a,1.7) -- (\b,0.5); \fill (\b,0.5) circle (1.6pt);}
  \node at (8.4,-0.15) {分支数 2、1、3：只能 \(2+1+3=6\)};
\end{tikzpicture}

\emph{乘法原理要求的不是每步面对同样的选项，而是每步面对同样多的选项。左树每个节点恰好分出两条，路径数才等于逐层相乘；右树分支参差，乘法无从下手，只能按第一层分类再相加。}
\end{center}

这个区分容易混淆，值得多说一句。从 \(n\) 个人里依次选出正副队长，第二步能选谁完全取决于第一步选了谁——两次的候选人集合根本不同。但两次的候选人个数是稳定的：第二步永远面对 \(n-1\) 个人。乘法要的是数目整齐，不是选项相同。回到开头那个网格：三个优化器对应的学习率个数分别是 2、4、4，数目不齐，乘法这才塌掉。

分支不齐时，退回加法。把结果集合切成互不相交的几块，各块分别计数再相加，这就是加法原理。它同样有一个必须守住的条件：各块不能重叠。一旦有对象同时落进两块，相加就把它数了两次；那种情形要留到本章第四节的容斥来收拾。

\subsection{排列、组合、重复字母，都是同一次除法}

真正把这几个公式串起来的，是一个比加法乘法更少被点名的原理。

\begin{proposition}[除法原理]
设 \(f:X\to Y\) 是满射，且每个 \(y\in Y\) 的原像恰好含 \(d\) 个元素，则 \(\abs Y=\abs X/d\)。
\end{proposition}

换成人话：如果目标集合里的每样东西，在源集合里都恰好有 \(d\) 个副本，那就先数源集合，再除以 \(d\)。难数的东西对应到好数的东西上去——这是整章反复用的动作，只不过这里的对应不是一一的，而是整齐的多对一。

\begin{center}
\begin{tikzpicture}[x=1cm,y=1cm,font=\footnotesize,>=stealth]
  \foreach \y/\t in {2.5/{甲乙丙},2.0/{甲丙乙},1.5/{乙甲丙},1.0/{乙丙甲},0.5/{丙甲乙},0.0/{丙乙甲}}{
    \node[left] at (0,\y) {\t};
    \draw[->,black!45] (0.25,\y) -- (4.0,1.25);}
  \draw[black!55,fill=black!5] (4.05,0.85) rectangle (6.4,1.65);
  \node at (5.2,1.25) {\{甲, 乙, 丙\}};
  \node[black!60] at (2.1,-0.55) {\(3!=6\) 个有序结果压成 1 个无序结果};
\end{tikzpicture}

\emph{排列到组合的这一步不是新公式，是同一张对应表：每个三人小组在有序世界里有 \(3!\) 个分身。数完有序的，除掉分身数，就得到无序的。}
\end{center}

有了这条，四个公式各归其位。全排列 \(n!\) 和依次取 \(k\) 个的排列数 \(P(n,k)=n!/(n-k)!\) 是纯乘法，逐步的候选人数为 \(n,n-1,\dots\)，整齐得很，不需要除。组合数
\[
\binom nk=\frac{P(n,k)}{k!}=\frac{n!}{k!\,(n-k)!}
\]
是一次除法：每个 \(k\) 人小组被它内部的 \(k!\) 种排列生成了 \(k!\) 遍。而 BANANA 这类含重复字母的排列，先当成六个互不相同的字母排出 \(6!\) 种，同一个可见结果被三个 A 的 \(3!\) 种内部交换和两个 N 的 \(2!\) 种内部交换重复生成了 \(3!\,2!\) 遍，于是答案是 \(6!/(3!\,2!)=60\)。一般地，总数 \(n\) 而各类各有 \(n_1,\dots,n_r\) 个的多重集排列数是 \(n!/(n_1!\cdots n_r!)\)。

三个公式，一次除法，分母写的都是同一件事：一个可见结果背后藏着多少个不可见的分身。忘了公式不要紧，问一句``我刚才把每个答案数了几遍''就能重新推出来。

\subsection{顺序算不算数，由问题决定，不由公式决定}

选还是排，从来不是符号问题。从 \(n\) 人中挑三人组成委员会，用 \(\binom n3\)；若三人分别担任主席、秘书、财务，交换他们就换出了一个新方案，用 \(P(n,3)\)。同一批人，同一个``选三个''的说法，答案差了 6 倍。

所以列式子之前先写一句话：一个结果由哪些信息完全确定？委员会的结果只由``哪三个人''确定，职务分工的结果还要加上``谁担任什么''。这句话写清楚了，顺序算不算、能不能重复、分支会不会重叠，全都随之确定。写不清楚就先别动笔——公式再熟也救不了一个没定义好的对象。

\subsection{组合爆炸是数据工作的常态背景}

计数在机器学习里最常出现的形态不是习题，而是``这个空间有多大''。\(d\) 个二值特征的所有子集有 \(2^d\) 个，\(d=50\) 时已超过 \(10^{15}\)；穷举式特征选择因此从来不是一个可执行的方案，能做的只有贪心、正则化或者别的近似。同样的数字也从反面起作用：一个模型能实现多少种不同的标注方式，正是刻画模型容量的出发点，样本量必须跟得上这个数的对数，否则模型有足够的自由度把训练集背下来。

值得注意的是量级。乘法原理下每加一个二值维度，空间就翻一倍；加法原理下每加一个分支，空间只增加一点。判断一个搜索能不能穷举，先看它的生成过程里有几层乘法。

\subsection{下一步：同一批对象数两遍}

本节的主张是，每个计数公式都该能由``对象怎样生成、被生成了几遍''重新推出来。这条主张还有一个更有力的用法：如果同一批对象可以按两种完全不同的方式生成，那么两种数法必须给出同一个数——不是应该，是必须。

下一节就把这条约束当成证明工具。它给出的不只是验算手段，还是一批组合恒等式的来源，以及一句在数据工作里天天用到的常识：同一个指标换一种口径统计，两个结果对不上，一定有一边错了。
